% Fragmento conservado; véase MANIFIESTO_ANTECEDENTES_LECTURA.json.
\section{Desarrollo del núcleo formal: condiciones iniciales, geometría y refinamiento}
\label{sec:extension}

El soporte de ochenta y una posiciones no coincide con el dominio de condiciones iniciales de la dinámica. Una condición inicial debe indicar además una orientación; y las hojas aditiva y multiplicativa disponen de cursores independientes. Esta distinción permite reconstruir exhaustivamente la emisión sin seleccionar de antemano una región numérica.

\subsection{Los estados iniciales y la recurrencia de emisión}

Para la enumeración utilizamos índices \(I_9=\{0,\ldots,8\}\), relacionados con las etiquetas positivas por \(i\mapsto i+1\). Sea
\[
 \mathcal S=I_9^2\times\{N,E,S,O\},\qquad
 \Omega_{\rm seed}=\mathcal S_+\times\mathcal S_\times.
\]
Cada condición inicial \(\omega=(s_+,s_\times)\) determina dos posiciones y dos direcciones. El dominio comprende todas las parejas ordenadas, por lo que
\begin{equation}
 |\mathcal S|=324,\qquad |\Omega_{\rm seed}|=324^2=104\,976.
 \label{eq:semillas}
\end{equation}
El anexo independiente enumera todas estas condiciones; el cuerpo presenta la ley que las transforma. Las filas de una tabla de emisiones son fibras de este dominio de condiciones iniciales.

La realización finita utiliza seis ventanas de nueve transiciones. En cada instante \(t\), la firma de fase es \(M_{\rm ph}(1+((t-1)\bmod9))\). Si vale \(+1\), se lee el residuo positivo \(\Sigma(i+1,j+1)=\rho_9(i+j+2)\) del primer cursor y después se desplaza ese cursor; si vale \(-1\), se lee el residuo positivo \(\Pi(i+1,j+1)=\rho_9((i+1)(j+1))\) del segundo cursor y después se desplaza el segundo. El valor cero incrementa el contador neutral. Al finalizar las transiciones veintisiete y cincuenta y cuatro, se invierten las orientaciones cardinales. Las evaluaciones enteras y sus cocientes se conservan en el registro, pero el emisor utiliza las lecturas residuales aquí especificadas.

En el sector de unidades, el lector multiplicativo utiliza el logaritmo discreto de base dos en \((\ZZ/9\ZZ)^\times\):
\[
 \ell_9(1)=0,\quad\ell_9(2)=1,\quad\ell_9(4)=2,\quad
 \ell_9(8)=3,\quad\ell_9(7)=4,\quad\ell_9(5)=5.
\]
Para la emisión finita se extiende con valor cero al sector no unitario. Esta extensión es un observable y no una identificación de los elementos \(3,6,9\) con la unidad multiplicativa. La hoja y su pertenencia al sector radical permanecen en el registro de trayectoria.

Sean \(S_m\) la suma de los tres residuos positivos aditivos \(\Sigma\) de la ventana \(m\), \(E_m\) la suma de los tres exponentes discretos de los residuos multiplicativos \(\Pi\) y \(Z_m=3\) el número de eventos neutrales. Así, \(S_m\) no suma las evaluaciones enteras \(S=i+j\), cuyos cocientes permanecen en el registro. Con \([a]_{10}\in\{0,\ldots,9\}\), la regla de emisión decimal es
\begin{align}
 d_{m,1}&=[S_m]_{10},\notag\\
 d_{m,2}&=[S_m+E_m]_{10},\notag\\
 d_{m,3}&=[3S_m+5E_m+7Z_m]_{10},\notag\\
 U_m&=100d_{m,1}+10d_{m,2}+d_{m,3}.
 \label{eq:emisor-seis}
\end{align}
Los coeficientes enteros y el calendario de esta recurrencia son parte explícita de la ley emisora declarada. No se reciben valores de \(\pi,\varphi,e\) ni de \(\alpha\). La distinción entre una regla de emisión y una constante reconocida a partir de su salida evita que la genealogía quede escondida en una tabla de decimales.

Si \(s_m=[S_m]_{10}\) y \(e_m=[E_m]_{10}\), la misma regla se escribe
\begin{equation}
 U_m=100s_m+10[s_m+e_m]_{10}+[3s_m+5e_m+1]_{10}.
 \label{eq:acoplamiento}
\end{equation}
El uno final es el residuo de \(7Z_m=21\) módulo diez. La reducción ternaria orientada del catálogo adopta la convención
\begin{equation}
 \Psi(U_1,\ldots,U_6)=([-U_1]_3,\ldots,[-U_6]_3).
 \label{eq:psi-catalogo}
\end{equation}
El signo de esta coordenatización es explícito: cambiarlo modifica las palabras y exige transportar las orientaciones posteriores. La identificación balanceada de cada componente de \(\FF_3\) con el TRIT se realiza después de \eqref{eq:psi-catalogo}.

\begin{proposition}[Factorización inyectiva de la emisión]
La palabra \(U=(U_1,\ldots,U_6)\) determina las dos firmas \(s=(s_1,\ldots,s_6)\) y \(e=(e_1,\ldots,e_6)\). La imagen de la emisión tiene \(18\cdot26=468\) elementos. La proyección ternaria tiene \(243\) valores distintos en el ambiente de \(729\) palabras.
\end{proposition}
\begin{proof}
Las centenas de \(U_m\) recuperan \(s_m\). Si \(b_m\) es su cifra de decenas, entonces \(e_m=[b_m-s_m]_{10}\). La cifra de unidades comprueba la tercera congruencia de \eqref{eq:acoplamiento}. Por ello el acoplamiento es inyectivo en el producto de firmas.

La enumeración de las \(324\) condiciones de un cursor mediante la recurrencia anterior produce dieciocho firmas aditivas, cada una con dieciocho preimágenes; y veintiséis multiplicativas, de las cuales veinticuatro tienen ocho preimágenes, una tiene veinticuatro y una tiene ciento ocho. El anexo conserva estas firmas y sus condiciones iniciales, permitiendo repetir la enumeración finita sin un catálogo numérico de referencia. La inyectividad del acoplamiento da \(18\cdot26\) emisiones y multiplicidades
\[
 \underbrace{432}_{18\cdot24}\text{ fibras de }144,
 \qquad18\text{ fibras de }432,
 \qquad18\text{ fibras de }1944.
\]
La suma es \(432\cdot144+18\cdot432+18\cdot1944=104\,976\). Aplicar \eqref{eq:psi-catalogo} a las \(468\) emisiones produce exactamente \(243\) palabras. Esta última cifra es el cardinal de la imagen, no el del espacio \(\FF_3^6\).
\end{proof}

Esta factorización explica la diferencia entre exhaustividad finita y profundidad ilimitada. El dominio \eqref{eq:semillas} es completo para el soporte y el emisor declarados; no contiene anticipadamente cada prefijo de cada prolongación. La profundidad pertenece al sistema de extensiones compatibles que se construye sobre esas condiciones.

\subsection{Clausura del calendario y avance de memoria}

El calendario de ciento ocho transiciones consta de doce ventanas nonádicas. Su estructura de grupo es más precisa que un producto informal de dos períodos:
\begin{equation}
 0\longrightarrow C_{12}\xrightarrow{\iota}C_{108}
 \xrightarrow{p}C_9\longrightarrow0,
 \quad\iota([b])=[9b],\quad p([t])=[t]_9.
 \label{eq:extension108}
\end{equation}
Si \(t=9b+r\) con \(0\le r<9\), la suma de dos fases produce el acarreo
\[
 (b,r)+(b',r')=
 \left(b+b'+\left\lfloor\frac{r+r'}9\right\rfloor,\ [r+r']_9\right),
\]
donde la primera coordenada se reduce módulo doce. La composición conserva el término que enlaza las dos coordenadas.

\begin{proposition}[No escisión del calendario]
La sucesión \eqref{eq:extension108} es exacta y no admite una sección que sea homomorfismo de grupos.
\end{proposition}
\begin{proof}
El núcleo de la reducción módulo nueve es el subgrupo de múltiplos de nueve, imagen de \(\iota\). Si la extensión se escindiera, por conmutatividad se tendría \(C_{108}\cong C_{12}\times C_9\). El primer grupo tiene exponente \(108\), mientras el segundo tiene exponente \(\operatorname{mcm}(12,9)=36\), contradicción.
\end{proof}

Nueve transiciones restauran la fase nonádica y avanzan una ventana; ciento ocho restauran el calendario finito. Ninguno de estos hechos obliga a borrar la trayectoria o a restaurar el prefijo de una prolongación. La holonomía nonádica \(\Gamma_9\) designa el transporte de una vuelta sobre estados con memoria, de modo que
\begin{equation}
 g(\Gamma_9x)=g(x),\qquad M(\Gamma_9x)=M(x)+1.
 \label{eq:holonomia-memoria}
\end{equation}
La igualdad de fase y el incremento de memoria son compatibles y constituyen dos observaciones del mismo transporte.


\subsection{Compresión observable y conservación ortogonal}

Una realización isométrica permite expresar de manera cuantitativa qué pierde una observación reducida. Sea \(U:\mathcal H\to\mathcal H\) una realización isométrica del transporte de una vuelta, \(U^*U=I\), y \(J:\mathcal H_{\rm vis}\to\mathcal H\) una inclusión isométrica. Definimos
\[
 T=J^*UJ,\qquad D=(I-JJ^*)UJ.
\]
Entonces
\begin{equation}
 UJ=JT+D,\qquad J^*D=0,\qquad I-T^*T=D^*D.
 \label{eq:conservacion-ortogonal}
\end{equation}
La última igualdad se deduce tomando el producto adjunto de la primera y usando la ortogonalidad. La contracción de la componente visible no constituye por sí sola destrucción de información del estado total.

\begin{proposition}[Conservación a horizonte finito]
Si \(I-T^*T=D^*D\), entonces, para todo \(n\ge1\),
\begin{equation}
 I=\sum_{r=0}^{n-1}T^{*r}D^*DT^r+T^{*n}T^n.
 \label{eq:telescopia}
\end{equation}
\end{proposition}
\begin{proof}
Multiplicar \(D^*D=I-T^*T\) por \(T^{*r}\) y \(T^r\) y sumar produce una suma telescópica. Se cancelan los términos interiores y queda \(I-T^{*n}T^n\).
\end{proof}

El último sumando conserva el estado terminal. Suprimirlo antes de justificar un límite alteraría la igualdad. Esta distinción es el contenido operativo de la conservación de la unidad en la compresión: la norma se distribuye entre la componente visible, la memoria y el terminal, sin ser reemplazada por un único residuo.


\subsection{Completación cúbica y dos nociones de frontera}

La ampliación del soporte cúbico de lado nueve al de lado diez admite una estratificación posicional exacta:
\begin{equation}
 10^3=9^3+3\cdot9^2+3\cdot9+1=729+243+27+1.
 \label{eq:cubo}
\end{equation}
El primer término corresponde a las tres coordenadas interiores; el segundo, a una coordenada nueva; el tercero, a dos; y el último, al vértice de tres coordenadas nuevas. La corona de ampliación contiene \(271\) posiciones. En cambio, la frontera interna del cubo discreto de lado nueve contiene \(9^3-7^3=386\). La diferencia entre estos recuentos es una diferencia de dominio, no un conflicto de valores.
